%%% ---------------------------------------------------------------------------
%%% setup.tex —— 论文元数据
%%%
%%% 只填内容，不写格式。所有排版规则在 acadarticle.cls 里。
%%% 对应 GB/T 7713.2—2022 前置部分的著录项。
%%% ---------------------------------------------------------------------------

%%% ---- 中文题名与作者 ----
%%% 题名一般不超过 20 字，不使用非公知的缩略语、字符、代号
\title{基于多尺度特征融合的遥感图像目标检测方法}

%%% 多作者用 \quad 分隔；上标数字对应下面的单位编号
\author{张三\textsuperscript{1}\quad 李四\textsuperscript{1,2}\quad 王五\textsuperscript{2}}

%%% 单位要写全称，到二级单位，附城市与邮编（GB/T 7713.2）
\affiliation{%
  1.~中国科学技术大学 信息科学技术学院，合肥 230027；\\
  2.~中国科学技术大学 大数据学院，合肥 230026%
}

%%% ---- 关键词：3~8 个，分号分隔，与英文一一对应 ----
\keywords{遥感图像；目标检测；多尺度特征；特征融合；卷积神经网络}

%%% ---- 中图分类号与文献标志码 ----
%%% 中图分类号查《中国图书馆分类法》；文献标志码 A=理论与应用研究论文
\clcnumber{TP391.4}
\doccode{A}

%%% ---- 英文题名、作者、单位（对应中文，供英文摘要块使用）----
\entitle{Multi-scale Feature Fusion for Object Detection in Remote Sensing Imagery}
\enauthor{ZHANG San\textsuperscript{1}, \quad LI Si\textsuperscript{1,2}, \quad WANG Wu\textsuperscript{2}}
\enaffiliation{%
  1.~School of Information Science and Technology, University of Science and
  Technology of China, Hefei 230027, China;\\
  2.~School of Data Science, University of Science and Technology of China,
  Hefei 230026, China%
}
\enkeywords{remote sensing imagery; object detection; multi-scale feature;
  feature fusion; convolutional neural network}

%%% ---- 基金项目与通信作者（自动排为首页无标记脚注；anonymous 模式下隐去）----
\funding{国家自然科学基金资助项目（62376001）；国家重点研发计划（2023YFB0000000）}
\corresponding{李四，教授，E-mail: lisi@ustc.edu.cn}

%%% ---- 页眉的短标题 ----
\runningtitle{基于多尺度特征融合的遥感图像目标检测方法}

%%% ---- 日期：留空则不显示 ----
\date{}

%%% ---- 参考文献数据库 ----
%%% refs.bib 是指向 ../common/refs.bib 的符号链接，与 beamer 模板共用同一份
%%% 文献库，改一处两处都生效。
\addbibresource{refs.bib}
